\section{Brechas tecnológicas y oportunidades de investigación}

A pesar de los avances tecnológicos recientes, existen importantes brechas que limitan el desarrollo de soluciones efectivas para la gestión de residuos orgánicos y la mitigación de emisiones de gases de efecto invernadero (GEI). Entre estas brechas destacan la falta de sistemas automatizados para monitorear variables ambientales críticas, como temperatura, humedad y pH, durante procesos biotecnológicos como la bioconversión mediante la mosca soldado negra (\textit{Hermetia illucens}) \parencite{kamilaris2017review}. Además, la precisión y confiabilidad de sensores IoT en entornos biológicos complejos sigue siendo un desafío técnico significativo. En paralelo, los modelos de aprendizaje automático (AAA) actuales carecen de capacidad para optimizar dinámicamente procesos no lineales, lo que restringe su aplicabilidad en contextos industriales o productivos.

Estas limitaciones tecnológicas tienen un impacto directo en la implementación de estrategias sostenibles para la gestión de residuos. La ausencia de herramientas robustas para Medición, Reporte y Validación (MRV) dificulta la trazabilidad y evaluación precisa de emisiones de GEI, así como la toma de decisiones informadas. Asimismo, la falta de sistemas integrados que combinen tecnologías emergentes con procesos biotecnológicos reduce la escalabilidad y replicabilidad de soluciones basadas en la Mosca Soldado Negra (\textit{Hermetia illucens}), limitando su potencial para contribuir a una economía circular más sostenible \parencite{surendraRethinkingOrganicWastes2020}.

No obstante, estas brechas tecnológicas también representan oportunidades claras para futuras investigaciones. El desarrollo de sensores IoT especializados y algoritmos de AAA capaces de analizar datos en tiempo real podría revolucionar la eficiencia operativa de procesos de bioconversión. Además, la creación de plataformas digitales para MRV permitiría mejorar la precisión en la medición de emisiones y facilitar la adopción de prácticas más sostenibles. Esta investigación propone abordar estas necesidades mediante el diseño e implementación de un sistema integral que combine tecnologías emergentes con procesos biotecnológicos avanzados, contribuyendo así a reducir costos, minimizar impactos ambientales y promover soluciones escalables para la gestión de residuos orgánicos \parencite{ahmedAgricultureClimateChange2020,vanIntegrationInternetofThingsSustainable2022}.